% =====================================================================
% EXTRACTION: the probability core of "Golden Reasoning" (master v125; rev 4, after the second external referee),
% prepared as a standalone paper for a stochastics audience, following
% the corpus's extraction convention (cf. nexic.tex).  Assembled by
% Fable 5, Aug 21 2026, from the master's Sec. gr1 and appendices, with
% the numerics of app:numerics RUN and reported in Sec. 6 (previously a
% specification only).  Program vocabulary neutralized per the
% extraction convention: "technostrength" -> the muffling exponent x;
% "technoanatomy" -> the response exponent rho; one remark records the
% companion names.  All theorem content is the master's; the abstract,
% introduction framing, Sec. 6 results, and discussion are new.
% CITATION NOTE: bibliography reconstructed from the master's keys;
% AUTHOR TO VERIFY every entry against the sources before circulation.
% =====================================================================
\documentclass[11pt]{article}
\usepackage[margin=1.1in]{geometry}
\usepackage{amsmath,amssymb,amsthm,mathtools}
\usepackage{booktabs}
\usepackage{graphicx}
\usepackage[colorlinks=true]{hyperref}
\usepackage{url}
\usepackage{enumitem}

\newtheorem{theorem}{Theorem}[section]
\newtheorem{proposition}[theorem]{Proposition}
\newtheorem{lemma}[theorem]{Lemma}
\newtheorem{corollary}[theorem]{Corollary}
\theoremstyle{definition}
\newtheorem{definition}[theorem]{Definition}
\newtheorem{ansatz}[theorem]{Ansatz}
\theoremstyle{remark}
\newtheorem{remark}[theorem]{Remark}

\newcommand{\PP}{\mathbb{P}}
\newcommand{\E}{\mathbb{E}}
\newcommand{\R}{\mathbb{R}}
\newcommand{\one}{\mathbf{1}}
\newcommand{\hm}{\mathsf{h}}
\newcommand{\Fc}{\mathcal{F}}
\newcommand{\Lc}{\mathcal{L}}
\newcommand{\Imm}{\mathsf{V}}
\DeclareMathOperator{\Leb}{Leb}

\title{Survivorship Conditioning for Hazard Rates Driven at the
Log-Derivative Level:\\ \large A Two-Channel Model, a Two-Condition
Viability Theorem, and the Critical Doob Limit}
\author{Ajax Benander\thanks{Email:
\href{mailto:akb9408@rit.edu}{akb9408@rit.edu}. This paper extracts, as
a self-contained development in applied probability, the hazard model
and survivorship theory of a larger framework; the companion
developments carry the program built on them and are cited, not
presupposed. Drafting collaboration: Anthropic Claude (Opus~4.x,
Fable~5), under the author's direction.}}
\date{September 2026 --- Working Draft, rev.~6}

\begin{document}
\maketitle

\begin{abstract}
We study the survival theory of a hazard rate whose noise is placed one
integration further out than is standard: rather than diffusing $h$ or
$\log h$, we diffuse the negative logarithmic derivative
$\mu=-\frac{d}{dt}\log h$, so that the integrated driver is a
Kolmogorov diffusion and the governing persistence exponent moves from
$\tfrac12$ to $\tfrac14$. The hazard carries two channels: an external
background, muffled by the accumulated exponent through a bounded
logistic factor, and an internal term $\Gamma e^{-\rho x}$ the same
exponent feeds or starves according to the sign of the response
exponent $\rho$. The paper's claim in one sentence: the
total hazard ever accrued stays finite exactly when the response
exponent is positive and the driver's long-run behavior is favorable,
so that no magnitude of accumulated capacity substitutes for the sign,
and no event of bounded duration on a surviving path
perfectly separates the survivorship-conditioned law from the
unconditioned one. Three results carry the claim. First, a
trichotomy in $\rho$: for $\rho\le0$ the viability event
$\{\Lambda_\infty<\infty\}$ is null for every parameter value and every
driver, with an exponential survival tail at frozen background, while
for $\rho>0$ viability factors exactly as an indicator in $\rho$ times
the driver's escape probability --- the classical scale-function
expression when a scale limit is finite, the stationary-mean
criterion in the positive-recurrent case ---
two independent necessary conditions, with a closed-form Gaussian
sigmoid in the canonical self-reinforcing case. Second, a transfer
theorem: the two-channel model's survival asymptotics coincide with
those of the bare integrated-driver model up to a shift of initial data
and bounded constants, conditional on a soft-wall reduction we state as
an explicit Ansatz and support numerically. Third, the survivorship
interpolation $\PP(\cdot\mid T>t)$ is a monotone flow of upward tilts
whose $t\to\infty$ limit, in the critical regime, is the pure Doob
$\hm$-transform by the persistence-harmonic function: locally
absolutely continuous with respect to the unconditioned law on
survivors, globally singular, approached at the polynomial rate $s/4t$
with the tail index as linear-response coefficient, and classified
against exponential-tail regimes into three modes. Monte Carlo
experiments with exact pair simulation confirm the $t^{-1/4}$ exponent
for the soft-killing functional, the boundedness (indeed apparent
convergence) of the soft-to-hard survival ratio, the gauge role of the
initial hazard, the finite-window kinematic factor $(1-s/t)^{-1/4}$,
self-similar convergence of the conditioned credit, and the
qualitative structure of the survivorship drift. One empirical
signature follows: under criticality the population hazard of
survivors decays as $1/(4t)$, mortality deceleration to zero at a
universal rate, against the plateau that static frailty selection
predicts.
\end{abstract}

\section{Introduction}\label{sec:intro}

Stochastic hazard rates are standard objects: the wrapper
$S_t=\exp(-\int_0^t h)$ with $h$ a diffusion is the reduced-form
framework of credit risk \cite{Lan04} and of stochastic mortality,
where the force of mortality is routinely modeled by a diffusion, the
mean-reverting Brownian Gompertz specification of Milevsky and
Promislow \cite{MP01} being a canonical instance. What distinguishes
the present model is not the wrapper but \emph{where the noise is
placed}: standard practice diffuses $h$ or $\log h$, whereas we diffuse
$\mu=-\tfrac{d}{dt}\log h$, one integration further out. The
consequences are structural rather than cosmetic: $\log h$ becomes
$C^1$, the state becomes a degenerate, $3/2$-self-similar hypoelliptic
diffusion, Kolmogorov's 1934 example \cite{Kol34,Hor67}, and the
governing persistence exponent moves from $\tfrac12$ to $\tfrac14$.

\subsection{The model in three lines}
\emph{The hazard.}
$h^{\mathrm{tot}}_t=B_t\varsigma_b(x_t)+\Gamma e^{-\rho x_t}$: an
external background $B_t$ attenuated by a bounded logistic muffling
factor $\varsigma_b$ that the accumulated exponent $x$ improves, plus
an internal term that $x$ feeds or starves according to the sign of
$\rho$.

\emph{The driver.} $x_t=\int_0^t\mu_s\,ds$ with $\mu$ Brownian: the
noise sits on $\mu=-\frac{d}{dt}\log h$. This placement is the paper's
principal modeling claim.

\emph{Departure.} $T$ is the first time the cumulative hazard exhausts
an independent unit exponential (a Cox construction), and viability is
the event $\{\Lambda_\infty<\infty\}$ that the total hazard ever
accrued stays finite. Departure with probability one and infinite life
expectancy are compatible, and both hold in the critical case.

\emph{Position in the branch.} This paper is the root of a four-paper
branch and imports nothing from the others: the theorem companion on
target-conditioned records cites it for the survival specialization,
the Route World companion for the conditioned law it hypothesizes
governs the record, and the technoanatomy companion for the viability
theorem and the ceiling; the hard-wall law $Q^{\mathrm{hard}}$ is
proved here and the soft law $Q^{\mathrm{soft}}$ is conditional on
Ansatz~\ref{ans:ratio}, and every downstream use says which it consumes.

\emph{The claim.} The total hazard ever accrued stays finite exactly when
the response exponent is positive and the driver's long-run behavior is
favorable, so that no magnitude of accumulated capacity substitutes for
the sign (Theorem~\ref{thm:twocond}), and no event of bounded duration on a
surviving path perfectly separates the survivorship-conditioned law
from the unconditioned one (Theorem~\ref{thm:singular}). The rest of the
paper establishes those two statements and the transfer that connects
the two-channel model to the classical persistence theory they rest
on.

\subsection{Relation to the literature}
Four bodies of work border this paper; what follows maps what is
classical, what is adjacent, and what is claimed as new.

\emph{The classical layer.} The state process is Kolmogorov's 1934
diffusion \cite{Kol34}, the primordial hypoelliptic example
\cite{Hor67}. Its persistence theory is a mature subject with a clean
lineage, McKean \cite{McK63}, Goldman \cite{Gol71}, Lachal
\cite{Lac91}, Sinai \cite{Sin92}, Isozaki--Watanabe \cite{IW94},
Denisov--Wachtel \cite{DW15}, surveyed in \cite{AS15} and, as the
\emph{random acceleration process}, possessing an extensive physics
literature reviewed in \cite{BMS13}. The conditioned (never-hitting)
process, our limit object in Section~\ref{sec:cond}, is constructed by
Groeneboom, Jongbloed and Wellner \cite{GJW99}. None of this is ours;
we import it.

\emph{The adjacent frameworks.} Our survivorship interpolation belongs
to two established genres. First, Theorem~\ref{thm:Q} is a
\emph{penalization} statement in the sense of Roynette--Vallois--Yor
\cite{RVY06}: a limit of Wiener-type measures reweighted by a
time-dependent functional, identified as a martingale change of
measure --- here the weight is survival itself, and the polynomial
tail places it in the clockless class. Second, the mode classification
of Theorem~\ref{thm:modes} sits against the theory of quasi-stationary
distributions and $Q$-processes surveyed in \cite{MV12}: classical QSD
theory is at home in the exponential-tail class, while the critical
polynomial mode admits no proper QSD in fixed coordinates, its Yaglom
object being self-similar. We use both frameworks as positioning, not
as inputs.

\emph{The conceptual ancestor.} That selection on survival makes the
observed ensemble diverge systematically from the unbiased one is the
frailty phenomenon of Vaupel, Manton and Stallard \cite{VMS79} and
Vaupel and Yashin \cite{VY85}: when the frail are removed first, the
population hazard falls below the individual hazard, producing
mortality deceleration at the population level with no corresponding
deceleration for any individual. We claim no priority over that
insight, and regard this paper as exhibiting its \emph{critical
dynamical} analog: frailty is a static random variable fixed at birth
and the analysis is a mixture computation, whereas here the randomness
is a driver evolving in time and the object is a limit of path
measures singular with respect to the unbiased law
(Theorem~\ref{thm:singular}).

\emph{What is claimed as new.} To our knowledge, the following are new
at least in assembly: the log-derivative placement of the noise, with
its gauge analysis and the two-sided soft-wall analysis whose
remaining gap is located exactly (Ansatz~\ref{ans:soft}); the response
trichotomy and the two-condition viability theorem
(Theorem~\ref{thm:twocond}); and the survivorship interpolation read
as a monotone flow with the tail index as its linear-response
coefficient, classified into saturation, transformation and
concentration (Theorem~\ref{thm:modes}). Section~\ref{sec:numerics}
reports Monte Carlo support for the one non-rigorous reduction.

\emph{Convention.} For positive functions, $f\sim g$ means $f/g\to1$;
$f\asymp g$ means $0<\liminf f/g\le\limsup f/g<\infty$. $\Phi$ is the
standard normal distribution function. All processes live on a
filtered space satisfying the usual conditions; $W$ is a standard
one-dimensional Brownian motion. Claims are tagged: [P] proved here,
[C] classical or cited, [A] conditional on Ansatz~\ref{ans:soft},
[S] scaling-level, [O] open.

\emph{Vocabulary.} Ten terms are the paper's own and are fixed here.
\emph{Departure}: the killing time $T$. \emph{Viability}: the event
$\{\Lambda_\infty<\infty\}$. \emph{Credit}: the accumulated exponent
$x_t=\int_0^t\mu$. \emph{Response exponent}: $\rho$, the sign of
which decides viability. \emph{Muffling factor}: the bounded logistic
$\varsigma_b(x)$ attenuating the background. \emph{Driver-viability factor}:
$V_{\mathrm{drv}}$, the second factor of Theorem~\ref{thm:twocond},
the driver's escape probability when it is transient and the indicator
of a positive stationary mean when it is positive recurrent. \emph{Hard wall / soft wall}: killing at
first hitting of a level, against killing by accrued hazard.
\emph{Survivorship interpolation}: the family
$\PP(\cdot\mid T>t)$. \emph{Bias field}: the drift the interpolation
manufactures. \emph{Modes}: saturation, transformation, concentration,
the three limits of Theorem~\ref{thm:modes}.

\section{The model}\label{sec:model}

\begin{definition}[Driftless model]\label{def:model}
Fix $\sigma>0$ and initial data $(\mu_0,D_0)\in\R\times\R$. Define
\begin{equation}\label{eq:model}
\mu_t=\mu_0+\sigma W_t,\qquad
D_t=D_0+\int_0^t\mu_s\,ds,\qquad
h_t=e^{-D_t},\qquad
\Lambda_t=\int_0^t h_u\,du .
\end{equation}
The credit $D$ and the hazard are one coordinate: the initial hazard
is $h_0:=e^{-D_0}$, and a multiplicative constant $q$ in the hazard
is the translation $D_0\mapsto D_0-\log q$ of the initial credit,
which is the content of Lemma~\ref{lem:gauge}. The hard wall of
Section~\ref{sec:persist} sits at $D=0$, where $e^{-D}$ passes
through one.
Let $E\sim\mathrm{Exp}(1)$ be independent of $W$ and
$T:=\inf\{t\ge0:\Lambda_t\ge E\}$, so that
$\PP(T>t\mid\Fc^W_\infty)=e^{-\Lambda_t}$. The \emph{viability event}
is $\Imm:=\{\Lambda_\infty<\infty\}$; on $\Imm$ the conditional
probability of immortality, $e^{-\Lambda_\infty}$, is strictly
positive, and
$\PP(T=\infty)=\E[e^{-\Lambda_\infty}\one_\Imm]\le\PP(\Imm)$.
\end{definition}

The pair $(\mu,D)$ is the Kolmogorov diffusion,
$\Lc=\tfrac{\sigma^2}{2}\partial_\mu^2+\mu\,\partial_D$, Gaussian with
covariance
$\bigl(\begin{smallmatrix}\sigma^2t&\sigma^2t^2/2\\
\sigma^2t^2/2&\sigma^2t^3/3\end{smallmatrix}\bigr)$
and hypoelliptic by H\"ormander's condition at first order
\cite{Kol34,Hor67}. Three structural facts are immediate and load
bearing. There is no boundary in $h$: the logarithmic parametrization
makes $h>0$ automatic, and the whole problem concerns the growth of
the integrated driver. The credit $D$ is nondecreasing exactly while
$\mu\ge0$: accumulated credit is spent only through a negative rate.
And the initial hazard is exactly gauge:

\begin{lemma}[Gauge role of the initial credit; {[P]}]\label{lem:gauge}
Writing $\Lambda^{(D_0)}$ for the cumulative hazard from initial
credit $D_0$, $\Lambda^{(D_0)}_t=e^{-D_0}\Lambda^{(0)}_t$ pathwise;
hence $\Imm$ does not depend on $D_0$, while $\PP(T=\infty)$ is
continuous and nondecreasing in $D_0$, with
$\lim_{D_0\uparrow\infty}\PP(T=\infty)=\PP(\Imm)$ and
$\lim_{D_0\downarrow-\infty}\PP(T=\infty)=0$. The initial credit,
equivalently the initial hazard scale $h_0=e^{-D_0}$, can never affect
a dichotomy, only a value.
\end{lemma}

\begin{remark}[Nonparametric-Bayes reading]
Definition~\ref{def:model} places an integrated Wiener prior on
$-\log h$ --- the classical cubic smoothing-spline prior of Wahba
\cite{Wah78} --- so the results below may be read as the survival
asymptotics that canonical smoothness prior implies. A stronger
reading is available: writing $\mathcal D[f]:=-\frac{d}{dt}\log f$,
one has $\mathcal D[\bar F]=h$ and $\mathcal D[h]=\mu$, so the driver
is the second iterate of the operation that defines the hazard, and
the model occupies a rung of a ladder rather than a member of a
family. Noise at level one ($\log h$ Brownian) gives persistence
exponent $\tfrac12$; at level two, the present model, $\tfrac14$; at
level three the credit is twice-integrated Brownian motion and the
exponent is open \cite{AS15,AD13}. Level two is the last rung that is
both nontrivial and solvable.
\end{remark}

\begin{definition}[Two-channel model]\label{def:twochannel}\footnote{The
two-channel form and the viability criterion $\Lambda_\infty<\infty$
have a dated germ in the author's document of August 2022
\cite{BenGolden22}, where the criterion appears as $e^{a/b}<\infty$
under a two-channel ``Outlook''; the present formulation replaces the
germ and is cited for priority only.}
Fix $\sigma,\epsilon>0$, $n\in\mathbb N$, $\Gamma>0$, $\rho\in\R$, and
offset $b>0$. With $\mu$, $x$ as in \eqref{eq:model} (writing $x$ for
the integrated driver), let
$B_t=\sum_{i=1}^n e^{\epsilon W^{(i)}_t}$ for Brownian motions
$W^{(i)}$ independent of one another and of the driver $W$, an
assumption the infinite-set lemma below requires when the set is
generated by the driver, for independent Brownian
motions $W^{(i)}$, and
\begin{equation}\label{eq:hazard}
h^{\mathrm{tot}}_t
=\underbrace{B_t\,\varsigma_b(x_t)}_{\text{external}}
+\underbrace{\Gamma\,e^{-\rho x_t}}_{\text{internal}},
\qquad
\varsigma_b(x):=\frac{1+e^{-b}}{1+e^{x-b}},
\end{equation}
normalized so that $\varsigma_b(0)=1$ exactly. Survival, $T$, and
$\Imm$ are as before with $h^{\mathrm{tot}}$ in place of $h$.
\end{definition}

\begin{proposition}[Channel dominance; boundedness; {[P]}]
\label{prop:dominance}
(i) For every $\rho>0$,
$\Gamma e^{-\rho x}/(B\varsigma_b(x))\to\infty$ as $x\to-\infty$, so
the background drops out of the total at large negative $x$ for every
positive $\rho$. (ii) As $x\to+\infty$ the external channel dominates
for all large $x$ if and only if $\rho>1$; for $\rho<1$ the two cross
at $x_\times=(k+b)/(1-\rho)$, $k:=\log(B/\Gamma)$. (iii)
$0<\varsigma_b(x)<1+e^{-b}$ for all $x$: the external channel can be
reduced without bound but amplified by at most $1+e^{-b}\to1$. The
internal channel is unbounded above.
\end{proposition}

\begin{remark}[The causal asymmetry]
Item (iii) enforces structurally a constraint a multiplicative form
$Be^{-x}$ would violate: a failure of the system's own muffling cannot
amplify the external world, while nothing bounds self-inflicted
hazard. In operator terms, at fixed background,
$\mathcal D[h^{\mathrm{int}}]=\rho\,\mu$ exactly, while
$\mathcal D[h^{\mathrm{ext}}]=\varsigma(x-b)\,\mu$ with logistic
weight: the internal channel preserves the defining operator relation
with constant gain $\rho$, and the external channel must be deformed
because boundedness is being purchased. This identifies $\rho$
intrinsically as the internal channel's operator gain,
$\rho=-\partial_x\log h^{\mathrm{int}}$: the rate at which accumulated
exponent converts into change in self-inflicted hazard. The companion
program names $x$ and $\rho$ \emph{technostrength} and
\emph{technoanatomy} \cite{GoldenMaster}; here they are the muffling
exponent and the response exponent.
\end{remark}

\section{The response trichotomy and the two-condition theorem}
\label{sec:trichotomy}

We now let $\rho$ range over $\R$.

\begin{proposition}[Trichotomy; {[P]}]\label{prop:trichotomy}
Write $h(x)=B\varsigma_b(x)+\Gamma e^{-\rho x}$ with $B,\Gamma>0$.
\begin{enumerate}[label=\emph{(\roman*)},itemsep=1pt]
\item If $\rho>0$, $h$ is strictly decreasing with $h(x)\to0$ as
$x\to+\infty$.
\item If $\rho=0$, $h$ is strictly decreasing with $h(x)\to\Gamma>0$:
protective against a floor.
\item If $\rho<0$, $h\to B(1+e^{-b})$ as $x\to-\infty$ and
$h\to\infty$ as $x\to+\infty$, so that $\inf_xh(x)>0$ for frozen
$B>0$, which is all the fatality theorem consumes. In the regime
$|\rho|>1$ there is a unique interior critical point, a minimum, whose location
is the large-$x$ asymptote of the exact critical-point equation
$Ky^{r-1}(1+y)^2=1$ below, obtained from $\varsigma_b(x)\sim
\mathrm{const}\cdot e^{-x}$ (writing $A:=Be^{b}$, $|\rho|=-\rho$):
\begin{equation}\label{eq:xstar}
x_\ast\;\sim\;\frac{\log\bigl(A/(|\rho|\Gamma)\bigr)}{1+|\rho|}
\qquad(x_\ast-b\to+\infty),
\qquad
h_\ast\;\sim\;A^{\frac{|\rho|}{1+|\rho|}}\,\Gamma^{\frac{1}{1+|\rho|}}
\Bigl(|\rho|^{\frac{1}{1+|\rho|}}+|\rho|^{-\frac{|\rho|}{1+|\rho|}}\Bigr).
\end{equation}
For $|\rho|\le1$ the number of interior critical points depends on the
parameters and may be zero: with $y=e^{x-b}$ and $r=|\rho|$ they solve
$Ky^{r-1}(1+y)^2=1$ for a positive constant $K$, whose left side has
logarithmic derivative $r-1+2y/(1+y)$, increasing throughout when
$r\ge1$ and decreasing then increasing when $r<1$, so that for $r<1$
there are zero, one tangential, or two solutions, and for $r=1$ none
when $K\ge1$. Nothing below depends on the count.
\end{enumerate}
\end{proposition}

\begin{lemma}[Infinite sets defeat the driftless background; {[P]}]
\label{lem:infbackground}
Let $W$ be a standard Brownian motion, $\epsilon>0$, and let
$S\subset[0,\infty)$ be a measurable set, possibly random but
independent of $W$, with $\mathrm{Leb}(S)=\infty$ almost surely. Then
$\int_S e^{\epsilon W_u}\,du=\infty$ almost surely; consequently the
same holds with $B_u=\sum_{i=1}^ne^{\epsilon W^{(i)}_u}\ge
e^{\epsilon W^{(1)}_u}$ in place of $e^{\epsilon W_u}$.
\end{lemma}

\begin{proof}
Conditioning on $S$, it suffices to treat deterministic $S$ of
infinite measure. For each $T$, finiteness of
$\int_{S\cap[T,\infty)}e^{\epsilon W_u}du$ is unchanged by the
multiplicative factor $e^{\epsilon W_T}$ --- it is decided by
$\int_{S\cap[T,\infty)}e^{\epsilon(W_u-W_T)}du$ --- and the integral
over $S\cap[0,T]$ is a.s.\ finite, so the event
$\{\int_Se^{\epsilon W_u}du=\infty\}$ lies in
$\sigma(W_u-W_T:u\ge T)$ for every $T$, hence in the tail
$\sigma$-field of the increments, and has probability $0$ or $1$ by
Kolmogorov's zero--one law. Were the probability $0$, then
$\int_Se^{\epsilon W_u}du<\infty$ a.s., and applying the same to $-W$,
also $\int_Se^{-\epsilon W_u}du<\infty$ a.s.; adding,
$2\,\mathrm{Leb}(S)\le\int_S(e^{\epsilon W_u}+e^{-\epsilon W_u})du
<\infty$, contradicting $\mathrm{Leb}(S)=\infty$.
\end{proof}

\begin{proposition}[Recurrent drivers; {[P]}]\label{prop:recurrent}
Suppose the driver of \S\ref{sec:model} is positive recurrent with
stationary law $\pi$ and $\int|\mu|\,\pi(d\mu)<\infty$; write
$\bar m:=\int\mu\,\pi(d\mu)$.
\begin{enumerate}[label=\emph{(\alph*)},itemsep=1pt]
\item If $\bar m>0$ then $D_t/t\to\bar m$ a.s., $\Lambda_\infty<\infty$
a.s., and $\PP(\Imm)=1$.
\item If $\bar m<0$ then $\Lambda_\infty=\infty$ a.s.\ and $T<\infty$ a.s.
\item If $\bar m=0$ and the cycle integral defined in the proof is
nondegenerate (automatic for nondegenerate diffusions), then
$\Lambda_\infty=\infty$ a.s.: \emph{the critical case falls on the mortal
side.}
\end{enumerate}
For null-recurrent drivers no general statement is offered
\textbf{[O]}; the Brownian driver, the fundamental null-recurrent case, is
settled unconditionally by the driftless analysis of \S\ref{sec:persist}.
\end{proposition}

\begin{proof}
(a), (b): the ergodic theorem for positive recurrent one-dimensional
diffusions gives $D_t/t\to\bar m$ a.s.\ (see \cite[Ch.~X]{RY99}); if
$\bar m>0$ then eventually $D_t\ge D_0+\tfrac{\bar m}2t$ and
$\Lambda_\infty<\infty$; if $\bar m<0$ the integrand
$e^{-(D_t-D_0)}\to\infty$.

(c) Fix interior levels $a<b$ and regenerate at up-crossings: let
$R_0:=\inf\{t:\mu_t=b\}$ and inductively
$S_k:=\inf\{t>R_{k-1}:\mu_t=a\}$, $R_k:=\inf\{t>S_k:\mu_t=b\}$. By the
strong Markov property the cycles over $[R_{k-1},R_k)$ are i.i.d.; positive
recurrence gives $\E[L]<\infty$ for the cycle length
$L:=R_1-R_0$, and the cycle occupation identity yields, for the cycle
integral $\xi_k:=\int_{R_{k-1}}^{R_k}\mu_s\,ds$,
\[
\E[\xi_1]=\E[L]\int\mu\,d\pi=0,\qquad
\E|\xi_1|\le\E\Bigl[\int_{\text{cycle}}|\mu_s|ds\Bigr]
=\E[L]\int|\mu|\,d\pi<\infty .
\]
Thus $S_n:=D_{R_n}-D_{R_0}=\sum_{k\le n}\xi_k$ is a nondegenerate mean-zero
i.i.d.\ walk, so $\liminf_nS_n=-\infty$ a.s.\ by the classical
oscillation dichotomy for random walks (see, e.g., \cite{Dur19});
in particular the events $B_n:=\{D_{R_n}\le-1\}$ occur infinitely often
a.s. Choose $\delta>0$ and $M<\infty$ with
\[
q:=\PP_b\Bigl(L\ge\delta,\ \sup_{0\le u\le\delta}
\Bigl|\int_0^u\mu_s\,ds\Bigr|\le M\Bigr)>0,
\]
possible since $L>0$ a.s.\ and the supremum is finite a.s. Let
$G_n$ be the corresponding event for cycle $n{+}1$; by the strong Markov
property $\PP(G_n\mid\Fc_{R_n})=q$, so
$\sum_n\PP(B_n\cap G_n\mid\Fc_{R_n})=q\sum_n\one_{B_n}=\infty$ a.s., and
L\'evy's extension of the Borel--Cantelli lemma (see, e.g.,
\cite{Dur19}) gives
$B_n\cap G_n$ infinitely often a.s. On $B_n\cap G_n$, for
$u\in[0,\delta]$ one has $D_{R_n+u}\le-1+M$, hence
\[
\int_{R_n}^{R_n+\delta}e^{-D_s}\,ds\ \ge\ \delta\,e^{\,1-M-D_0}\,>\,0 ,
\]
and these windows are disjoint. Summing over infinitely many of them,
$\Lambda_\infty=\infty$ a.s.
\end{proof}

\begin{theorem}[Non-positive response is fatal, and fast; {[P]}]
\label{thm:trapped}
Let $\rho\le0$. Then: (i) $\Lambda_\infty=\infty$ almost surely, for
the random background and regardless of the driver --- pathwise in the
driver, almost surely in the background. (ii) If the background is
frozen, $B_t\equiv B>0$, then $h_\ast=\inf_xh(x)>0$ and
$\PP(T>t)\le e^{-h_\ast t}$, whence $\E[T]<\infty$. For the random
background no positive floor exists ($\inf_tB_t=0$ a.s.), so the
exponential rate is frozen-background (annealed) only, and whether an
unconditional rate holds is open [O]. (iii) No power-law structure
survives.
\end{theorem}

% CORRECTION RECORD (external referee flow-through, Aug 22 2026):
% (i) was stated "along every continuation" and its proof asserted
% background divergence "along every continuous path", which is false
% path-by-path; the claim is a.s.\ and now routes through
% Lemma~lem:infbackground.  (ii)'s old random-background clause
% conditioned on the null event {inf B >= B0 > 0} and was vacuous;
% withdrawn, with the annealed scoping made explicit.

\begin{proof}
(i) If $x_t\to+\infty$ the internal term diverges pointwise (for
$\rho<0$; for $\rho=0$ it is the constant $\Gamma$). If
$x_t\to-\infty$ the external term tends to $B_t(1+e^{-b})$ and
$\int_0^\infty B_t\,dt=\infty$ a.s.\ by
Lemma~\ref{lem:infbackground} with $S=[0,\infty)$. On any other
driver path, either $\mathrm{Leb}\{t:x_t\ge0\}=\infty$, where the
trough floors the hazard at $\Gamma e^{-\rho x_t}\ge\Gamma$
deterministically, or $\mathrm{Leb}\{t:x_t\le0\}=\infty$, where the
crest with $S=\{t:x_t\le0\}$ routes through
Lemma~\ref{lem:infbackground}. (ii) $\Lambda_t\ge h_\ast t$ and
$\PP(T>t)=\E[e^{-\Lambda_t}]\le e^{-h_\ast t}$.
\end{proof}

\begin{theorem}[Two conditions for viability; {[P]}]\label{thm:twocond}
Let the driver be a regular diffusion on $\R$ with scale function $s$,
and $\rho\in\R$. Define the \emph{driver-viability functional}
\begin{equation}\label{eq:vdrv}
V_{\mathrm{drv}}(\mu_0)=
\begin{cases}
\PP_{\mu_0}(\mu_t\to+\infty)=\dfrac{s(\mu_0)-s(-\infty)}{s(+\infty)-s(-\infty)},
& \text{transient driver,}\\[8pt]
1, & \text{positive recurrent, }\bar m>0,\\[2pt]
0, & \text{positive recurrent, }\bar m\le0,\\[2pt]
0, & \text{Brownian driver,}\\[2pt]
\text{open [O]}, & \text{general null recurrence,}
\end{cases}
\end{equation}
with $\bar m=\int\mu\,d\pi$ the stationary mean where it exists, the
transient formula read with its boundary conventions ($s(+\infty)=\infty$
and $s(-\infty)$ finite giving $0$; $s$ bounded above with
$s(-\infty)=-\infty$ giving $1$). Then
\begin{equation}\label{eq:twocond}
\PP\bigl(\Lambda_\infty<\infty\bigr)
=\underbrace{\one\{\rho>0\}}_{\text{response}}
\cdot
\underbrace{V_{\mathrm{drv}}(\mu_0)}_{\text{driver}}.
\end{equation}
The two factors are logically distinct necessary conditions, not
independent events: the first is a sign, the second is silent on it.
In the positive recurrent case the driver does not escape; its time
integral acquires positive linear drift, which is why the functional
is stated piecewise rather than as a single scale-function ratio.
\end{theorem}

% CORRECTION RECORD (external referee flow-through, Aug 22 2026):
% the previous clause read the second factor as 0 whenever
% s(+inf)=inf; an Ornstein--Uhlenbeck driver mean-reverting to a
% positive level (both scale limits infinite, stationary mean > 0)
% gives viability 1 by the ergodic theorem.  The counterexample is
% the referee's; the restatement routes the doubly-infinite case
% through Proposition~prop:recurrent.

\begin{proof}
On $\{\mu_t\to+\infty\}$: eventually $\mu_t\ge1$, so
$x_t\ge x_{t_0}+(t-t_0)$ and, for $\rho>0$,
$h(x_t)\lesssim e^{-\min(1,\rho)(t-t_0)}$, whence
$\Lambda_\infty<\infty$. For $\rho\le0$,
Theorem~\ref{thm:trapped}(i) makes the first factor necessary. On
$\{\mu_t\to-\infty\}$, $x_t\to-\infty$ and
$h\to B(1+e^{-b})>0$, so $\Lambda_\infty=\infty$; on the recurrent
event the outcome is Proposition~\ref{prop:recurrent} ---
$\bar m>0$ gives $D_t/t\to\bar m$ and $\Lambda_\infty<\infty$,
$\bar m\le0$ (nondegenerate) gives $\Lambda_\infty=\infty$ --- with
the Brownian driver settled by the driftless analysis and general
null recurrence left open, which is the statement's scope clause.
The hitting probabilities are the classical scale-function formula
\cite{KT81,RY99}.
\end{proof}

\begin{corollary}[The solvable case; {[P]}]\label{cor:solvable}
For the self-reinforcing driver $d\mu_t=\nu\mu_t\,dt+\sigma\,dW_t$,
$\nu>0$ --- the member of the driver family whose drift is
proportional to its state ---
$s'(v)=e^{-\nu v^2/\sigma^2}$ is integrable at both ends and
\begin{equation}\label{eq:sigmoid}
\PP\bigl(\Lambda_\infty<\infty\bigr)
=\one\{\rho>0\}\cdot\Phi\!\Bigl(\frac{\mu_0\sqrt{2\nu}}{\sigma}\Bigr):
\end{equation}
a sharp threshold in the response exponent multiplied by a smooth
Gaussian sigmoid of width $\sigma/\sqrt{2\nu}$ in the initial rate.
With $\rho<0$, feedback accelerates failure:
$\mu_t\to+\infty$ drives $h\to\infty$ superexponentially, so a system
with adverse response and feedback is strictly worse off than one with
adverse response and none.
\end{corollary}

\begin{remark}[Fluctuating response; {[S]}]\label{rem:fluct}
If $\rho=\rho_t$ is a process, an adverse excursion of duration
$\delta$ at exponent level $x$ contributes
$\asymp\delta\,\Gamma e^{|\rho|x}$ to $\Lambda$, so survivability
requires $\delta\lesssim\Gamma^{-1}e^{-|\rho|x}$: the affordable
duration of adverse response decays exponentially in the accumulated
exponent, and a $\rho_t$ recurrent on $\R$ forces
$\Lambda_\infty=\infty$ whatever the driver. Viability then requires
$\rho_t$ eventually positive almost surely, and \eqref{eq:twocond}
becomes a product of two scale functions. Supplying dynamics for
$\rho_t$ is the model's principal open direction [O].
\end{remark}

\begin{proposition}[Perturbation in $\Gamma$, and the threshold;
fixed-time expansion {[P]}, long-time threshold {[O]}]\label{prop:perturb}
Writing $Q_{\mathrm{ext}}$ for the external-only survivorship measure,
\begin{equation}\label{eq:perturb}
\PP(T>t)=\PP_{\mathrm{ext}}(T>t)\Bigl[1-\Gamma\,
\E_{Q_{\mathrm{ext}}}\!\Bigl[\int_0^te^{-\rho x_s}ds\Bigr]
+O(\Gamma^{2})\Bigr],
\end{equation}
the fixed-$t$ expansion being a Taylor--Feynman--Kac identity. Under
$Q_{\mathrm{ext}}$ the credit grows as $x_s\asymp\sigma s^{3/2}$ on
the scaling picture, which suggests that the first-order coefficient
converges as $t\to\infty$ if and only if $\rho>0$; the suggestion is
not a proof, since $\E_Q[e^{-\rho x_s}]$ is controlled by the mass the
conditioned process keeps near low credit and not by its typical
scale, and the long-time threshold is entered as open. Were it
established, the expansion's radius would collapse exactly at the
trichotomy's boundary, the signature of a genuine transition.
\end{proposition}

\begin{theorem}[Adverse credit occupies half of logarithmic time in the
critical regime; {[P]}]\label{thm:half}
For the driftless driver, almost surely
\begin{equation}\label{eq:half}
\lim_{S\to\infty}\frac1S\,\Leb\bigl\{s\in[0,S]:D_{e^{s}}-D_0<0\bigr\}=\frac12,
\end{equation}
so the set of times at which the credit sits below its initial value
has logarithmic density exactly one half, for every initial datum. The
statement is confined to the critical recurrent regime and fails for
any driver whose credit is transient.
\end{theorem}

\begin{proof}
$Z_s:=(D_{e^{s}}-D_0)/(\sigma e^{3s/2})$ is a stationary, centered,
Gaussian, ergodic process, by the Lamperti transform of the integrated
Brownian motion under its $3/2$-self-similarity, so Birkhoff's ergodic
theorem gives logarithmic density $\PP(Z_0<0)=\tfrac12$ by symmetry.
\end{proof}

\noindent The theorem is exported for the technoanatomy companion, which
reads it as the statement that adverse strength is a generic condition
in the critical regime and not evidence of adverse anatomy.

\section{Persistence and the transfer}\label{sec:persist}

\begin{theorem}[Persistence of the hard-wall functional; {[C]}]
\label{thm:persist}
Let $\tau:=\inf\{t\ge0:D_t\le0\}$ for the driftless model started at
$\mu_0\ge0$, $D_0>0$. Then
$\PP_{(\mu,d)}(\tau>t)\sim c\,\hm(\mu,d)\,t^{-1/4}$ as $t\to\infty$,
where $\hm$ is the persistence-harmonic function of the killed
Kolmogorov semigroup, with boundary behavior
$\hm(\mu,0^+)\propto\mu^{1/2}$ and $\hm(0^+,d)\propto d^{1/6}$, and
$\E_x[\hm(X_s)\one_{\{\tau>s\}}]=\hm(x)$. The exponent is Sinai's
\cite{Sin92}; the refined form and the harmonic function are
McKean, Goldman, Lachal, Isozaki--Watanabe, and
Groeneboom--Jongbloed--Wellner
\cite{McK63,Gol71,Lac91,IW94,GJW99}; universality is
Denisov--Wachtel \cite{DW15}, surveyed in \cite{AS15,BMS13}.
\end{theorem}

The soft-killing functional replaces the hard wall by the killing rate
$e^{-D}$, exponentially negligible for $D\gg0$ and overwhelming for
$D\ll0$: an absorbing wall of thickness $O(1)$ in $D$ at the level
$D=0$ where the hard wall $\tau=\inf\{t:D_t\le0\}$ sits, negligible
against excursions of size $t^{3/2}$. Killing approximates absorption
when the credit crosses downward through zero, and the sign convention
is fixed here once so that it cannot drift. The reduction is stated as two hypotheses of increasing strength,
because the exponent and the conditioned law require different
things: the tail asymptotics rest on the weaker and the $Q$-process
on the stronger.

\begin{ansatz}[A: soft killing preserves the exponent]\label{ans:soft}
For every admissible initial datum there exist
$0<c\le C<\infty$ and $t_0$ with
$c\,\PP(\tau>t)\le\PP(T>t)\le C\,\PP(\tau>t)$ for $t\ge t_0$, after an
$O(1)$ adjustment of the initial credit. Equivalently
$\log\PP(T>t)=\log\PP(\tau>t)+O(1)$.
\end{ansatz}

\begin{ansatz}[B: soft survival has a ratio asymptotic]\label{ans:ratio}
Writing $K_tf(x)=\E_x[e^{-\int_0^tq(X_u)du}f(X_t)]$ for the
Feynman--Kac semigroup with $q(\mu,D)=e^{-D}$, there is a positive
function $g$ on the state space and a constant $c_{\mathrm{soft}}$ with
$K_t1(x)\sim c_{\mathrm{soft}}\,g(x)\,t^{-1/4}$ locally uniformly in
$x$; formally $g$ is the zero-energy positive solution of
$(\Lc-q)g=0$ with the growth conditions the soft wall imposes.
\end{ansatz}

\noindent Ansatz~\ref{ans:soft} does not imply Ansatz~\ref{ans:ratio}:
comparability fixes $\log\PP_x(T>t)$ to $O(1)$ but leaves the hidden
constant free to depend on $x$, so that $\PP_x(T>t)\sim c(x)\hm(x)t^{-1/4}$
with $c$ non-constant is compatible with A and incompatible with the
conclusion that the soft conditioned law is the hard-wall
$\hm$-transform. The question the two hypotheses leave open, which
the numerical agenda of Section~\ref{sec:numerics} targets, is
whether $g=C\,\hm(\cdot+\delta)$ for some effective wall translation
$\delta$.

One half is rigorous modulo a cited input: on the event
$\{D_u\ge2\log\frac{2+u}{2}\ \forall u\le t\}$ the cumulative hazard
is bounded, and a logarithmic moving boundary does not change the
persistence exponent of an integrated process \cite{AS15,DW15}; the
upper half reduces to the plausible, unproved statement that
persistence with a bounded occupation tolerance has the same exponent
as strict persistence. The nearest standing literature cuts against
assuming this lightly: persistence with \emph{partial survival}, each crossing killing with probability $p$, carries exponents
varying continuously with $p$ \cite{MB98}, for the random acceleration
process among others \cite{Bur14}, so softening a wall can shift an
exponent. What distinguishes the present softening is its structure:
killing is rate-based and confined to a boundary layer of width $O(1)$
in the credit against excursions of scale $t^{3/2}$, a vanishing
fraction of the range, where partial survival taxes every crossing
at fixed price at every scale. Everything tagged [A] depends on exactly
this gap, and Section~\ref{sec:numerics} tests it directly; the
measured constancy of the soft-to-hard ratio is evidence the
boundary-layer case sits in the exponent-preserving class the
moving-boundary results delimit from the rigorous side.

\begin{proposition}[Transfer; scaling statement {[S]}, prefactor
conjectural]\label{thm:transfer}
In the two-channel model with $\rho>1$ the persistence exponent
$1/4$ is expected to survive the two-channel perturbation, because
the $t^{3/2}$ scale of the integrated driver dominates the $t^{1/2}$
scale of the background's fluctuations, so that under
Ansatz~\ref{ans:soft} $\log\PP(T>t)=-\tfrac14\log t+O(1)$. A heuristic derivation, not reproduced here and not relied on, suggests
the more specific form
\begin{equation}\label{eq:transfer}
\PP(T>t)\;\asymp\;
\hm\bigl(\mu_0-\tfrac{\epsilon^2}{2},\,D_0-\log n-b\bigr)\cdot
e^{-\Theta(n(b/\sigma)^{2/3})}\cdot t^{-1/4},
\end{equation}
every difference between the two-channel model and the driftless one
appearing as a shift of initial data or a bounded multiplicative
constant; that form is entered here as a conjecture, its shifts and
prefactor not established by the hypotheses stated in this paper, and
the numerics below test the bare soft wall rather than the full
two-channel process. Where the exponent holds, $\E[T]=\infty$ while
$T<\infty$ a.s., residual life is asymptotically Pareto$(\tfrac14)$,
and quantiles grow as the fourth power of rarity.
\end{proposition}

\section{Conditioning on survival}\label{sec:cond}

\begin{definition}[Survivorship interpolation; bias field]
\label{def:interp}
$\PP^{(t)}:=\PP(\cdot\mid T>t)$, likewise $\PP^{(t)}_\tau$ for the
hard wall; for $0\le s<t$ the bias field on $[0,s]$ is
\begin{equation}\label{eq:bias}
\beta_t(s):=\frac{d\PP^{(t)}}{d\PP}\Big|_{\Fc_s}
=\one_{\{T>s\}}\,\frac{\PP_{X_s}(T>t-s)}{\PP_{X_0}(T>t)},
\qquad X_s=(\mu_s,D_s).
\end{equation}
\end{definition}

\begin{lemma}[Survivorship tilts upward; {[P]}]\label{lem:monotone}
Coupling two initial data componentwise with the same $W$ shows both
$\PP_x(T>t)$ and $\PP_x(\tau>t)$ nondecreasing in $(\mu_0,D_0)$;
hence every bias field is, on the survival event, a nondecreasing
function of the state: the interpolation is a monotone flow of upward
tilts in rate and credit.
\end{lemma}

\begin{theorem}[The critical limit is the pure Doob transform;
{[C]}, soft version {[A]}]\label{thm:Q}
Under the refined asymptotics of Theorem~\ref{thm:persist}, for every
$s\ge0$ and bounded $\Fc_s$-measurable $\Psi$,
$\E[\Psi\mid\tau>t]\to\E_Q[\Psi]$ as $t\to\infty$, where $Q$ is the
time-homogeneous Markov law with
\begin{equation}\label{eq:Qdensity}
\frac{dQ}{d\PP}\Big|_{\Fc_s}
=\one_{\{\tau>s\}}\,\frac{\hm(X_s)}{\hm(X_0)}:
\end{equation}
the pure Doob $\hm$-transform at eigenvalue one. Under
Ansatz~\ref{ans:ratio}, and conditionally on it, the soft
survivorship law $\PP(\cdot\mid T>t)$ converges to the Feynman--Kac
Doob transform
\begin{equation}\label{eq:Qsoft}
\frac{dQ^{\mathrm{soft}}}{d\PP}\Big|_{\Fc_s}
=e^{-\int_0^sq(X_u)du}\,\frac{g(X_s)}{g(X_0)},
\end{equation}
which coincides with \eqref{eq:Qdensity} exactly when $g$ is
proportional to $\hm$ after a credit translation. Ansatz~\ref{ans:soft}
alone yields neither the limit nor its existence.
\end{theorem}

\begin{proof}
With $f_t:=\one_{\{\tau>s\}}\PP_{X_s}(\tau>t-s)/\PP_{X_0}(\tau>t)$,
the Markov property gives $\E[f_t]=1$; factorizing,
$f_t=\one_{\{\tau>s\}}
\frac{(t-s)^{1/4}\PP_{X_s}(\tau>t-s)}{t^{1/4}\PP_{X_0}(\tau>t)}
(1-s/t)^{-1/4}\to f:=\one_{\{\tau>s\}}\hm(X_s)/\hm(X_0)$ a.s.; the
martingale property gives $\E[f]=1$, so $f_t\to f$ in $L^1$ by
Scheff\'e, and $\E[\Psi f_t]\to\E_Q[\Psi]$.
\end{proof}

\begin{remark}[Why survivorship has no clock here; {[P]}]
\label{rem:noclock}
For regularly varying tails the kinematic factor
$\PP_x(\tau>t-s)/\PP_x(\tau>t)\to1$ for fixed $s$: the age of the
conditioning drops out and the limit is an un-tilted $h$-transform.
For exponential tails the same factor converges to $e^{\theta s}$ and
survivorship acquires a clock. In the penalization taxonomy
\cite{RVY06} the survival weight lands in the clockless class
precisely because its expectation is regularly varying. Quantitatively
\begin{equation}\label{eq:rate}
\beta_t(s)=\one_{\{\tau>s\}}\frac{\hm(X_s)}{\hm(X_0)}
\Bigl(1-\frac st\Bigr)^{-1/4}(1+o(1)),
\qquad
\Bigl(1-\frac st\Bigr)^{-1/4}=1+\frac{s}{4t}+O\Bigl(\frac{s^2}{t^2}\Bigr):
\end{equation}
the persistence exponent is the linear-response coefficient of
survivorship bias.
\end{remark}

\begin{theorem}[Survivorship drift; promotion to transience;
{[P]} given {[C]} inputs]\label{thm:drift}
Under $Q$ the state solves
$d\mu_t=\sigma^2\partial_\mu\log\hm(\mu_t,D_t)\,dt+\sigma\,dW^Q_t$,
$dD_t=\mu_t\,dt$: conditioning adds drift only in the noisy
coordinate, the drift is everywhere nonnegative, equals
$\sigma^2/2\mu$ near zero credit (the Bessel-$2$ form, the halving of
the classical exponent under one integration), and vanishes as
$\kappa:=\sigma^2D/\mu^3\to\infty$. Under $\PP$ the credit is oscillatory,
$\liminf D_s=-\infty$ and $\limsup D_s=+\infty$ (recurrence
classifications belong to the Markov pair, and $D$ alone is not
Markov); under $Q$ it is transient, $D_s\to\infty$ a.s.: the
P\'olya $2\mathrm D\to3\mathrm D$ transition with no literal change of
dimension anywhere --- survivorship alters the recurrence class of the
ensemble, not the state dynamics.
\end{theorem}

\begin{theorem}[Locally ours, globally elsewhere; {[P/C]}]
\label{thm:singular}
$Q(\tau=\infty)=1$; $\hm(X_s)\to\infty$ and $D_s\to\infty$ $Q$-a.s.
\cite{GJW99}; $\liminf_sD_s=-\infty$ $\PP$-a.s. Hence
$Q|_{\Fc_s}\ll\PP|_{\Fc_s}$ for every finite $s$, with the explicit
density \eqref{eq:Qdensity} and mutual absolute continuity on the
survival event, while $Q\perp\PP$ on $\Fc_\infty$: the survivor's
law is locally equivalent to the unconditioned one on survivors and
globally disjoint from it, the exit from the measure class completed
only at $t=\infty$. Consequently no event of bounded duration perfectly separates the
conditioned law from the unconditioned one: none has probability one
under one law and zero under the other on their common survival
support. Mutually absolutely continuous laws can still be far apart
statistically, so this is a statement about certification and not
about discrimination.
\end{theorem}

\begin{theorem}[Three canonical survivorship regimes; {[P/C/S]}]
\label{thm:modes}
Three tail classes of the driver produce three regimes of the
interpolation, canonical rather than exhaustive: stretched-exponential
and other subexponential tails, and regularly varying tails of
exponent other than $1/4$, fall under the same categories at the
level of the survival weight's behavior in the penalization taxonomy
but are not treated here. \emph{(I) Saturation} [P] (favorable:
$\PP(T=\infty)>0$): $\|\PP^{(t)}-\PP(\cdot\mid
T=\infty)\|_{\mathrm{TV}}\to0$ with every bias bounded by
$1/\PP(T=\infty)$; no new measure class. \emph{(II) Transformation}
[C/A] (regularly varying tails that also admit a statewise ratio
asymptotic, $\PP_x(T>t)\sim g(x)L(t)t^{-\alpha}$; regular variation of
the scalar tail alone does not determine the state-dependent ratio):
the clockless Doob or Feynman--Kac transform by $g$, locally absolutely
continuous, globally singular, with its kinematic factor; the hard-wall
case $\alpha=1/4$, $g=\hm$, is the proved instance
(Theorem~\ref{thm:Q}) and the soft case is conditional on
Ansatz~\ref{ans:ratio}. \emph{(III) Concentration} [S]
(exponential tails, $\PP(T>t)\sim e^{-\theta t}$): on fixed windows
the eigen-tilted transform, the quasi-stationary regime with a clock;
for the linearly adverse driver $d\mu=\nu\,dt+\sigma\,dW$ with
$\nu<0$ the clock is $\theta=3\nu^2/8\sigma^2$ and, at macroscopic
scale, the conditioned path law collapses onto the deterministic
optimal profile $\mu_{tv}/t\to\tfrac{|\nu|}4v(2-3v)$, a law of large
numbers for the rare, a large-deviation computation for the linearly
adverse driver that is stated here at scaling level and not derived. Modes (II) and (III) share $\PP(\Imm)=0$ and are separated by
the interpolation --- by whether the survivor's measure class is
approached polynomially, without a clock, or exponentially, with one.
\end{theorem}

\section{Numerical experiments}\label{sec:numerics}

The one load-bearing heuristic, Ansatz~\ref{ans:soft}, admits a direct
test, and three falsifiers are fixed before any result is read: a
fitted soft-wall slope departing robustly from $-1/4$; a dependence of
that slope on the initial hazard $h_0$; and instability of the
conditioned rescaled credit across late windows. Any one firing would
count against the Ansatz. All experiments use the exact
pair simulation: for a step $\Delta$, the increments
$(\xi,\eta)=(W_{t+\Delta}-W_t,\int_t^{t+\Delta}(W_s-W_t)ds)$ are
centered Gaussian with covariance
$\bigl(\begin{smallmatrix}\Delta&\Delta^2/2\\
\Delta^2/2&\Delta^3/3\end{smallmatrix}\bigr)$, so the update
$\mu\mathrel{+}=\sigma\xi$, $D\mathrel{+}=\mu\Delta+\sigma\eta$
reproduces the law at grid points with zero discretization error; the
only approximations are the trapezoidal quadrature of $\Lambda$ and
discrete monitoring of the wall. Parameters: $\sigma=1$,
$\mu_0=0$, with the soft-wall runs labeled by the initial hazard
$h_0=e^{-D_0}$ (so $h_0\in\{0.3,1,3,30\}$ is $D_0\in\{1.20,0,-1.10,-3.40\}$,
the gauge of Lemma~\ref{lem:gauge} realized as a credit translation);
$N=1.5\times10^5$ paths; a uniform grid of step
$0.025$ to $t=40$ (the zone where the $\Lambda$-mass accrues),
geometric thereafter at ratio $1.01$ to $t=3\times10^4$; slopes fitted
by least squares on $\log_{10}\PP$ against $\log_{10}t$ over
$t\in[10^2,3\times10^4]$, with path-bootstrap confidence intervals
($150$ resamples). A control run at half the step sizes
($0.0125$, ratio $1.005$) reproduces the estimates below to within
their intervals, excluding discretization bias.

\begin{table}[h]\centering\small
\begin{tabular}{@{}llll@{}}
\toprule
quantity & measured & predicted & status\\
\midrule
hard-wall slope, $\log\PP(\tau>t)$ & $-0.2499\ [-0.254,-0.243]$ &
$-1/4$ (theorem) & validates code\\
soft-wall slope, $h_0=0.3$ & $-0.2487$ & $-1/4$ [A] & supports\\
soft-wall slope, $h_0=1$ & $-0.2480\ [-0.255,-0.242]$ & $-1/4$ [A] &
confirms\\
soft-wall slope, $h_0=3$ & $-0.2483$ & $-1/4$ [A] & supports\\
soft-wall slope, $h_0=30$, late window & $-0.240$ & $-1/4$;
transient delayed & confirms gauge (exact)\\
ratio $\PP(T>t)/\PP(\tau>t)$, $h_0=1$, $t\ge10^2$ &
$\in[0.1768,0.1788]$ & bounded [A] & supports, sharply\\
kinematic coefficient (grid-exact) & $0.2874\pm0.0133$ & $0.2859$ &
confirms \eqref{eq:rate}\\
KS distance, $D_t/t^{3/2}$ at $t=2\times10^3$ vs $3\times10^4$ &
$0.010$ & $\to0$ (Yaglom) & consistent\\
\bottomrule
\end{tabular}
\caption{\small Monte Carlo results ($N=1.5\times10^5$; a second seed
reproduces the slope estimates to $\pm0.004$). The kinematic test
evaluates $(\PP(\tau>t{-}s)/\PP(\tau>t)-1)(t/s)$ at the grid spacing
$s/t=0.197$, whose exact power-law value is
$((1{-}s/t)^{-1/4}{-}1)(t/s)=0.2859$; the measurement matches the full
factor, not merely its linearization $1/4$.}
\label{tab:numerics}
\end{table}

Three findings deserve emphasis. First, the soft-killing functional
carries the $t^{-1/4}$ exponent at every tested $h_0$, and the
$h_0=30$ case exhibits exactly the behavior Lemma~\ref{lem:gauge}
predicts: a steepened transient over early windows (fitted slope
$-0.298$ over $t\ge10^2$) relaxing toward $-1/4$ as the window moves
late --- the gauge acts on the timescale, never the exponent. Second,
and this is the sharpest support for Ansatz~\ref{ans:soft}, the ratio
$\PP(T>t)/\PP(\tau>t)$ is not merely bounded over the fitted two
decades but constant to three decimal places
(Figure~\ref{fig:ansatz}b), consistent with convergence to a limit
constant $\approx0.178$ at these initial data --- a stronger behavior
than the Ansatz asserts. Third, the survivorship drift's qualitative
structure is confirmed: among paths surviving to the horizon, the
conditional increment $\E[\Delta\mu]$ regressed on $1/\mu$ over a
window $\Delta t=50$ at $t=2\times10^3$ gives a positive coefficient
of the predicted order at small $\kappa$ ($51\pm27$ against the
instantaneous-drift value $25$; the finite window biases upward, as
conditioned near-wall paths must gain rate over it) and a coefficient
indistinguishable from zero at $\kappa>2$ ($0.1\pm0.4$), matching
$G(\kappa)\to0$. The remaining sharp test, the prefactor constant
in \eqref{eq:transfer} and the deep-tail behavior at
$t/t_\ast\gtrsim10^{12}$ via multilevel splitting or
harmonic-function importance sampling, is left to a dedicated
computation. The experiments test Ansatz~\ref{ans:soft}, the exponent; they do
not test Ansatz~\ref{ans:ratio}. The computation that would is fixed
here: estimate $R_t(x,y)=\PP_x(T>t)/\PP_y(T>t)$ over a grid of initial
states $(\mu,D)$ at increasing $t$, and ask whether it converges to
$\hm(x)/\hm(y)$, or to $\hm(x+\delta)/\hm(y+\delta)$ for a fixed
translation, or to a ratio of some other positive $g$; with
harmonic-function importance sampling the deep tail is reachable.
\emph{None of the three falsifiers fired}: no fitted slope departs robustly from
$-1/4$, no dependence of the slope on $h_0$ appears, and the
conditioned rescaled credit is stable.

\begin{figure}[h]\centering
\includegraphics[width=0.95\textwidth]{fig_ansatz.pdf}
\caption{\small (a) Hard- and soft-wall survival curves against the
$t^{-1/4}$ reference. (b) The soft-to-hard ratio at $h_0=1$ over the
fitted window: flat to three decimals across two decades.}
\label{fig:ansatz}
\end{figure}

\section{Discussion}\label{sec:discussion}

The claim stated at the outset stands as proved: viability holds
exactly when the response exponent is positive and the driver's
long-run behavior is favorable, no magnitude of accumulated capacity
substituting for the sign, and no event of bounded duration on a surviving
path perfectly separates the conditioned law from the unconditioned
one. One
empirical signature follows. Under criticality the
population hazard of survivors is $-\frac{d}{dt}\log\PP(T>t)\sim
\frac{1}{4t}$: mortality deceleration not to a plateau but to zero, at
a universal rate set by the persistence exponent. The
reliability-theory literature on late-life mortality deceleration
\cite{GG01} documents plateaus as the fingerprint of static frailty
selection; a hazard decaying as $c/t$ with $c$ universal would be the
corresponding fingerprint of the dynamical mechanism, and
distinguishing the two in longevity data is a concrete empirical
question this model makes precise.

The theory assembled here is the critical dynamical analog of the
frailty phenomenon \cite{VMS79,VY85}: where frailty computes how an
observed hazard curve bends when a static susceptibility is selected
on, the present interpolation identifies the limiting law of the
survivor's entire trajectory, together with the drift conditioning
manufactures, and locates the observed-versus-unbiased divergence in
the measure class itself. The two-condition theorem gives the model's
structural moral in one line: an index of accumulated capacity tracks
the driver-viability factor and is silent on the sign that multiplies it, so
no such index can distinguish the ascent that arrives from the ascent
that is a fall --- and with $\rho<0$, feedback strictly accelerates
the fall. The interpretation of these results within a larger program, where $x$ and $\rho$ are read as a biosphere's capability and its
disposition to turn capability on itself, and the conditioning as
survivorship over deep time, is developed in the companion
\cite{GoldenMaster}, on which nothing above depends.

Five problems are open, and they are the paper's every [O] tag
gathered in one place. Supplying dynamics for $\rho_t$
(Remark~\ref{rem:fluct}) would convert the trichotomy's sign test into
a two-transience condition and is the model's principal open
direction. Closing Ansatz~\ref{ans:soft}, persistence with bounded
occupation tolerance, would discharge every [A] tag; the numerics
above make the outcome unsurprising without making it a theorem. The
deep-tail computation with harmonic-function importance sampling
would supply the one constant this paper leaves unevaluated. The
non-positive-response theorem's exponential rate is established at
frozen background only, and whether an unconditional rate holds is
open. And the recurrent-driver proposition and the two-condition
theorem are settled for the Brownian driver, the fundamental
null-recurrent case, and open for null-recurrent drivers in
general.

\begin{thebibliography}{99}\small
% AUTHOR TO VERIFY all entries against the sources before circulation.
\bibitem{Kol34} A.~Kolmogorov, Zuf\"allige Bewegungen (zur Theorie der
Brownschen Bewegung), \emph{Ann. of Math.} \textbf{35} (1934), 116--117.
\bibitem{Hor67} L.~H\"ormander, Hypoelliptic second order differential
equations, \emph{Acta Math.} \textbf{119} (1967), 147--171.
\bibitem{McK63} H.~P. McKean, A winding problem for a resonator driven
by a white noise, \emph{J. Math. Kyoto Univ.} \textbf{2} (1963),
227--235.
\bibitem{Gol71} M.~Goldman, On the first passage of the integrated
Wiener process, \emph{Ann. Math. Statist.} \textbf{42} (1971),
2150--2155.
\bibitem{Lac91} A.~Lachal, Sur le premier instant de passage de
l'int\'egrale du mouvement brownien, \emph{Ann. Inst. H. Poincar\'e
Probab. Statist.} \textbf{27} (1991), 385--405.
\bibitem{Sin92} Ya.~G. Sinai, Distribution of some functionals of the
integral of a random walk, \emph{Theoret. Math. Phys.} \textbf{90}
(1992), 219--241.
\bibitem{IW94} Y.~Isozaki and S.~Watanabe, An asymptotic formula for
the Kolmogorov diffusion and a refinement of Sinai's estimates for the
integral of Brownian motion, \emph{Proc. Japan Acad. Ser. A}
\textbf{70} (1994), 271--276.
\bibitem{DW15} D.~Denisov and V.~Wachtel, Random walks in cones,
\emph{Ann. Probab.} \textbf{43} (2015), 992--1044.
\bibitem{AS15} F.~Aurzada and T.~Simon, Persistence probabilities and
exponents, in \emph{L\'evy Matters V}, Lecture Notes in Math. 2149,
Springer, 2015, 183--224.
\bibitem{AD13} F.~Aurzada and S.~Dereich, Universality of the
asymptotics of the one-sided exit problem for integrated processes,
\emph{Ann. Inst. H. Poincar\'e Probab. Statist.} \textbf{49} (2013),
236--251.
\bibitem{BMS13} A.~J. Bray, S.~N. Majumdar, and G.~Schehr, Persistence
and first-passage properties in nonequilibrium systems, \emph{Adv.
Phys.} \textbf{62} (2013), 225--361.
\bibitem{GJW99} P.~Groeneboom, G.~Jongbloed, and J.~A. Wellner,
Integrated Brownian motion, conditioned to be positive, \emph{Ann.
Probab.} \textbf{27} (1999), 1283--1303.
\bibitem{RVY06} B.~Roynette, P.~Vallois, and M.~Yor, Some
penalisations of the Wiener measure, \emph{Jpn. J. Math.} \textbf{1}
(2006), 263--290.
\bibitem{MV12} S.~M\'el\'eard and D.~Villemonais, Quasi-stationary
distributions and population processes, \emph{Probab. Surv.}
\textbf{9} (2012), 340--410.
\bibitem{Lan04} D.~Lando, \emph{Credit Risk Modeling: Theory and
Applications}, Princeton Univ. Press, 2004.
\bibitem{MP01} M.~A. Milevsky and S.~D. Promislow, Mortality
derivatives and the option to annuitise, \emph{Insur. Math. Econ.}
\textbf{29} (2001), 299--318.
\bibitem{VMS79} J.~W. Vaupel, K.~G. Manton, and E.~Stallard, The
impact of heterogeneity in individual frailty on the dynamics of
mortality, \emph{Demography} \textbf{16} (1979), 439--454.
\bibitem{VY85} J.~W. Vaupel and A.~I. Yashin, Heterogeneity's ruses:
some surprising effects of selection on population dynamics,
\emph{Amer. Statist.} \textbf{39} (1985), 176--185.
\bibitem{Wah78} G.~Wahba, Improper priors, spline smoothing and the
problem of guarding against model errors in regression, \emph{J. Roy.
Statist. Soc. Ser. B} \textbf{40} (1978), 364--372.
\bibitem{KT81} S.~Karlin and H.~M. Taylor, \emph{A Second Course in
Stochastic Processes}, Academic Press, 1981.
\bibitem{RY99} D.~Revuz and M.~Yor, \emph{Continuous Martingales and
Brownian Motion}, 3rd ed., Springer, 1999.
\bibitem{Dur19} R.~Durrett, \emph{Probability: Theory and
Examples}, 5th ed., Cambridge University Press, 2019.
\bibitem{MB98} S.~N. Majumdar and A.~J. Bray, Persistence with partial
survival, \emph{Phys. Rev. Lett.} \textbf{81} (1998), 2626--2629.
\bibitem{Bur14} T.~W. Burkhardt, First passage of a randomly
accelerated particle, in \emph{First-Passage Phenomena and Their
Applications}, World Scientific, 2014, 21--44.
\bibitem{GG01} L.~A. Gavrilov and N.~S. Gavrilova, The reliability
theory of aging and longevity, \emph{J. Theoret. Biol.} \textbf{213}
(2001), 527--545.
\bibitem{BenGolden22} A.~Benander, Golden Science, August 2022,
\url{https://aethus.org/aethic-code/files/GoldenScience_Nov2022.pdf}.
\bibitem{GoldenMaster} A.~Benander, Golden Reasoning: soft horizons,
survivorship, and the two-channel muffled hazard, master manuscript,
2026.
\end{thebibliography}

\end{document}